\chapter{Ergodic theorems}
\label{chap:ergodic}

This chapter develops the three classical ergodic theorems that drive the
multiplicative ergodic theorem: the \emph{maximal ergodic inequality} of
Hopf and Garsia, the \emph{pointwise (Birkhoff) ergodic theorem}, and
\emph{Kingman's subadditive ergodic theorem}. The first is the analytic
gateway; the second turns the gateway into almost-everywhere convergence of
Birkhoff averages; the third is the genuine engine of Oseledets, turning a
subadditive sequence of log-norms into an almost-everywhere limit.

Throughout, \(T : X \to X\) is a map on a measurable space carrying a measure
\(\mu\), \texttt{birkhoffSum}\,\(T\,g\,n\,x = \sum_{k<n} g(T^{[k]}x)\) is the
\(n\)-th Birkhoff partial sum, and \texttt{birkhoffAverage}\,\(\R\,T\,g\,n\,x =
n^{-1}\,\texttt{birkhoffSum}\,T\,g\,n\,x\) is the Birkhoff average. The
arguments are due to Garsia, Katznelson--Weiss, and Karlsson; the formalization
keeps every quantity in \(\R\) (or, where \(-\infty\) cannot be excluded a
priori, in \texttt{EReal}) to avoid junk values.

\section{Maximal ergodic inequality}

\begin{definition}[Garsia's maximal function]\label{maxBirkhoff}
  \lean{Oseledets.maxBirkhoff}\leanok
  For a measurable \(T\), a function \(g : X \to \R\), \(N : \N\) and \(x\),
  the \emph{maximal function} is
  \[
    \texttt{maxBirkhoff}\,T\,g\,N\,x
      = \max_{0 \le k \le N} \texttt{birkhoffSum}\,T\,g\,k\,x,
  \]
  the nonempty \texttt{Finset.sup'} of the Birkhoff partial sums over
  \(k \in \texttt{range}\,(N+1)\). Since the \(k = 0\) term equals
  \(\texttt{birkhoffSum}\,T\,g\,0\,x = 0\), the maximal function is its own
  positive part.
\end{definition}

\begin{lemma}[Nonnegativity of the maximal function]\label{maxBirkhoff_nonneg}
  \lean{Oseledets.maxBirkhoff_nonneg}\leanok
  For all \(g, N, x\) one has \(0 \le \texttt{maxBirkhoff}\,T\,g\,N\,x\).
\end{lemma}
\begin{proof}\leanok
  \uses{maxBirkhoff}
  The index \(k = 0\) lies in \(\texttt{range}\,(N+1)\), and the corresponding
  term is \(\texttt{birkhoffSum}\,T\,g\,0\,x = 0\). The supremum over a
  nonempty finite set dominates each of its terms, so the maximal function is
  at least \(0\).
\end{proof}

\begin{lemma}[Recursion for the maximal function]\label{maxBirkhoff_succ}
  \lean{Oseledets.maxBirkhoff_succ}\leanok
  For all \(g, N, x\),
  \[
    \texttt{maxBirkhoff}\,T\,g\,(N+1)\,x
      = \texttt{birkhoffSum}\,T\,g\,(N+1)\,x
        \sqcup \texttt{maxBirkhoff}\,T\,g\,N\,x.
  \]
\end{lemma}
\begin{proof}\leanok
  \uses{maxBirkhoff}
  This is the identity \(\texttt{range}\,(N+2) = \texttt{insert}\,(N+1)\,
  (\texttt{range}\,(N+1))\) applied to the supremum: splitting off the top
  index gives the join of the new term with the previous maximum. Both
  inequalities follow from the universal/existential characterizations of
  \texttt{Finset.sup'}. The recursion is what makes integrability provable by
  induction on \(N\).
\end{proof}

\begin{lemma}[Garsia's pointwise inequality]\label{maxBirkhoff_le_add}
  \lean{Oseledets.maxBirkhoff_le_add}\leanok
  On the set where the maximal function is positive, for every \(x\) with
  \(0 < \texttt{maxBirkhoff}\,T\,g\,N\,x\),
  \[
    \texttt{maxBirkhoff}\,T\,g\,N\,x
      \le g\,x + \texttt{maxBirkhoff}\,T\,g\,N\,(T x).
  \]
\end{lemma}
\begin{proof}\leanok
  \uses{maxBirkhoff,maxBirkhoff_nonneg}
  Pull the constant \(g\,x\) through the supremum: by the additive shift of a
  nonempty \texttt{sup'} and the recursion
  \(\texttt{birkhoffSum}\,T\,g\,(k+1)\,x = g\,x +
  \texttt{birkhoffSum}\,T\,g\,k\,(Tx)\), the right-hand side equals
  \(\max_{0 \le k \le N} \texttt{birkhoffSum}\,T\,g\,(k+1)\,x\), the maximum of
  the \emph{shifted} partial sums. The maximum defining the left-hand side is
  attained at some index \(k_0\); positivity forces \(k_0 \ge 1\) (the index
  \(0\) gives the value \(0\)), so \(\texttt{birkhoffSum}\,T\,g\,k_0\,x\) is one
  of the shifted sums and is therefore dominated by their maximum.
\end{proof}

\begin{lemma}[Measurability and integrability of the maximal function]\label{integrable_maxBirkhoff}
  \lean{Oseledets.integrable_maxBirkhoff}\leanok
  If \(T\) is measure-preserving and \(g\) is integrable, then each Birkhoff
  partial sum and the maximal function \(\texttt{maxBirkhoff}\,T\,g\,N\) are
  integrable; if moreover \(g\) and \(T\) are measurable, then
  \(\texttt{maxBirkhoff}\,T\,g\,N\) is measurable.
\end{lemma}
\begin{proof}\leanok
  \uses{maxBirkhoff,maxBirkhoff_succ}
  Each summand \(g \circ T^{[j]}\) is integrable because \(T^{[j]}\) is
  measure-preserving, so the finite Birkhoff sum is integrable; measurability
  is likewise a finite sum of measurable compositions. Integrability of the
  maximal function follows by induction on \(N\) using the recursion
  \cref{maxBirkhoff_succ} and the fact that a join of two integrable functions
  is integrable; measurability is the measurability of a finite \texttt{sup'}.
\end{proof}

\begin{proposition}[Garsia's inequality at a fixed level]\label{setIntegral_maxBirkhoff_pos_nonneg}
  \lean{Oseledets.setIntegral_maxBirkhoff_pos_nonneg}\leanok
  For a measure-preserving \(T\) and a measurable integrable \(g\), and every
  \(N : \N\),
  \[
    0 \le \int_{\{x \,:\, 0 < \texttt{maxBirkhoff}\,T\,g\,N\,x\}} g \, d\mu.
  \]
\end{proposition}
\begin{proof}\leanok
  \uses{maxBirkhoff_le_add,maxBirkhoff_nonneg,integrable_maxBirkhoff}
  Write \(E\) for the level set and \(M = \texttt{maxBirkhoff}\,T\,g\,N\).
  Integrating the pointwise inequality \cref{maxBirkhoff_le_add} over \(E\)
  gives \(\int_E (M - M \circ T) \le \int_E g\). Now \(\int_E M \circ T \le
  \int_X M \circ T = \int_X M\) by nonnegativity and measure-preservation (via
  \texttt{integral\_map} and \(\texttt{map}\,T\,\mu = \mu\), avoiding any
  embedding hypothesis), and \(\int_X M = \int_E M\) since \(M\) vanishes off
  \(E\). Chaining these cancels the two maximal-function integrals and leaves
  \(0 \le \int_E (M - M\circ T) \le \int_E g\).
\end{proof}

\begin{lemma}[Level sets exhaust the target set]\label{iUnion_setOf_maxBirkhoff_pos}
  \lean{Oseledets.iUnion_setOf_maxBirkhoff_pos}\leanok
  The level sets \(\{x : 0 < \texttt{maxBirkhoff}\,T\,g\,N\,x\}\) are monotone
  in \(N\), and their union over all \(N\) is exactly
  \(\{x : \exists n,\ 0 < \texttt{birkhoffSum}\,T\,g\,(n+1)\,x\}\).
\end{lemma}
\begin{proof}\leanok
  \uses{maxBirkhoff}
  Monotonicity is \texttt{Finset.sup'} monotonicity over the nested ranges.
  For the union: \(0 < \texttt{maxBirkhoff}\,T\,g\,N\,x\) means some
  \(\texttt{birkhoffSum}\,T\,g\,k\,x > 0\) with \(k \le N\); the index \(0\)
  gives \(0\) and is excluded, so \(k = n+1\) for some \(n\). Conversely a
  positive partial sum at index \(n+1\) makes the maximal function positive at
  level \(N = n+1\).
\end{proof}

\begin{theorem}[Maximal ergodic inequality (Hopf--Garsia)]\label{setIntegral_birkhoffSum_pos_nonneg}
  \lean{Oseledets.setIntegral_birkhoffSum_pos_nonneg}\leanok
  For a measure-preserving \(T\) and an integrable \(f : X \to \R\),
  \[
    0 \le \int_{\{x \,:\, \exists n,\ 0 < \texttt{birkhoffSum}\,T\,f\,(n+1)\,x\}}
      f \, d\mu.
  \]
\end{theorem}
\begin{proof}\leanok
  \uses{setIntegral_maxBirkhoff_pos_nonneg,iUnion_setOf_maxBirkhoff_pos}
  First treat a measurable integrable \(g\): by
  \cref{setIntegral_maxBirkhoff_pos_nonneg} each level-set integral is
  nonnegative, and these integrals converge to the integral over the union
  \cref{iUnion_setOf_maxBirkhoff_pos} by monotone convergence of set integrals,
  so the limit is nonnegative. For general integrable \(f\), pass to a
  measurable representative \(g =^{\text{a.e.}} f\); the Birkhoff sums of \(f\)
  and \(g\) agree a.e., so the two target sets agree a.e.\ (the one for \(g\)
  being measurable, the one for \(f\) null-measurable), and the integrands
  agree a.e., which transfers the inequality from \(g\) to \(f\).
\end{proof}

\section{Birkhoff}

\begin{theorem}[Conditional expectation commutes with the dynamics]\label{condExp_invariants_comp_self}
  \lean{Oseledets.condExp_invariants_comp_self}\leanok
  For a finite measure, a measure-preserving measurable \(T\) and integrable
  \(g\), the conditional expectation onto the \(\sigma\)-algebra
  \(\texttt{invariants}\,T\) of \(T\)-invariant sets is a.e.\ \(T\)-invariant:
  \[
    \big(\mu[g \mid \texttt{invariants}\,T]\big) \circ T
      =^{\text{a.e.}} \mu[g \mid \texttt{invariants}\,T].
  \]
\end{theorem}
\begin{proof}\leanok
  Set-integral invariance of \(h \circ T\) over a measurable invariant set
  reduces, via \(\texttt{map}\,T\,\mu = \mu\), to the integral of \(h\). Hence
  \(\mu[g \circ T \mid \mathcal I]\) and \(\mu[g \mid \mathcal I]\) have equal
  integrals over every invariant set, so by uniqueness of the conditional
  expectation they agree a.e.; combining this with the commutation identity
  \(\mu[g\circ T \mid \mathcal I] =^{\text{a.e.}} (\mu[g\mid\mathcal I])\circ T\)
  (again by uniqueness, using that \(T\) is \((\mathcal I,\mathcal I)\)-measurable)
  yields the invariance.
\end{proof}

\begin{lemma}[Subexponential orbital tail]\label{ae_tendsto_orbit_div_atTop_zero}
  \lean{Oseledets.ae_tendsto_orbit_div_atTop_zero}\leanok
  For a measure-preserving \(T\) and integrable \(g\), the orbital tail
  \(n^{-1}\, g(T^{[n]}x)\) tends to \(0\) for \(\mu\)-a.e.\ \(x\).
\end{lemma}
\begin{proof}\leanok
  A Borel--Cantelli argument. For each threshold \(\delta = 1/(k+1)\) the series
  \(\sum_n \mu\{x : (n+1)\delta \le \abs{g\,x}\}\) is finite: the pointwise count
  of crossed thresholds is at most \(\abs{g\,x}/\delta\), and integrating
  (Tonelli) bounds the series by \(\delta^{-1}\int \abs{g}\). Measure-preservation
  transfers finiteness to the shifted orbit \(g \circ T^{[n]}\), so a.e.\ only
  finitely many \(n\) cross any fixed threshold; choosing \(k\) large makes the
  tail eventually smaller than any \(\varepsilon\).
\end{proof}

\begin{lemma}[A.e.\ boundedness of Birkhoff averages]\label{ae_bddAbove_birkhoffAverage}
  \lean{Oseledets.ae_bddAbove_birkhoffAverage}\leanok
  For a finite measure, a measure-preserving \(T\) and integrable \(g\), the
  Birkhoff averages \(n \mapsto \texttt{birkhoffAverage}\,\R\,T\,g\,(n+1)\,x\)
  are a.e.\ bounded above (and, applied to \(-g\), a.e.\ bounded below).
\end{lemma}
\begin{proof}\leanok
  \uses{setIntegral_birkhoffSum_pos_nonneg}
  The maximal ergodic inequality applied to \(g - c\) yields, for the maximal
  set \(B_c = \{x : \exists n,\ c < \texttt{birkhoffAverage}\,\R\,T\,g\,(n+1)\,x\}\),
  the estimate \(c\,\mu B_c \le \int_{B_c} g \le \int \abs{g}\). Taking
  \(c = k \in \N\) gives \(\mu B_k \le \int\abs{g}/k \to 0\), so the
  intersection \(\bigcap_k B_k\) is null. Off this null set the range of
  Birkhoff averages is bounded above.
\end{proof}

\begin{lemma}[Limsup invariance and vanishing perturbations]\label{limsup_birkhoffAverage_comp_ae}
  \leanok
  The pointwise limsup \(x \mapsto \limsup_n
  \texttt{birkhoffAverage}\,\R\,T\,g\,n\,x\) is a.e.\ \(T\)-invariant.
\end{lemma}
\begin{proof}\leanok
  \uses{ae_tendsto_orbit_div_atTop_zero,ae_bddAbove_birkhoffAverage}
  The difference \(A_n(g)(Tx) - A_n(g)(x) = n^{-1}(g(T^{[n]}x) - g\,x)\) tends
  to \(0\) a.e.\ by the tail estimate \cref{ae_tendsto_orbit_div_atTop_zero}.
  Two bounded sequences differing by a null sequence have equal limsup (proved
  via \texttt{limsup\_add\_const} and an \(\varepsilon\)-argument), and
  boundedness holds a.e.\ at both \(x\) and \(Tx\) by
  \cref{ae_bddAbove_birkhoffAverage}; hence the limsup is a.e.\ \(T\)-invariant.
\end{proof}

\begin{proposition}[The core maximal-inequality step]\label{measure_setOf_lt_limsup_eq_zero}
  \lean{Oseledets.measure_setOf_lt_limsup_eq_zero}\leanok
  For a finite measure, measure-preserving \(T\), integrable \(g\) and
  \(\varepsilon > 0\), the superlevel set where the limsup of the Birkhoff
  averages exceeds \(\mu[g\mid\mathcal I] + \varepsilon\) is null.
\end{proposition}
\begin{proof}\leanok
  \uses{setIntegral_birkhoffSum_pos_nonneg,condExp_invariants_comp_self,limsup_birkhoffAverage_comp_ae}
  Write \(L = \mu[g\mid\mathcal I]\) and \(Ls = \limsup A_\bullet(g)\). The set
  \(E = \{L + \varepsilon < Ls\}\) is a.e.\ \(T\)-invariant
  (\cref{condExp_invariants_comp_self}, \cref{limsup_birkhoffAverage_comp_ae}),
  hence a.e.\ equal to a genuinely invariant measurable set \(E'\). Feed
  \(\varphi = \mathbf 1_{E'}(g - L - \varepsilon)\) to the maximal ergodic
  inequality \cref{setIntegral_birkhoffSum_pos_nonneg}; on \(E'\) the partial
  sums of \(\varphi\) telescope (using orbit-constancy of \(L\)), and the
  maximal set equals \(E'\). Thus
  \(0 \le \int_{E'}(g - L - \varepsilon) = -\varepsilon\,\mu E'\) (using
  \(\int_{E'} g = \int_{E'} L\)), forcing \(\mu E' = 0\).
\end{proof}

\begin{theorem}[Pointwise (Birkhoff) ergodic theorem]\label{tendsto_birkhoffAverage_ae}
  \lean{Oseledets.tendsto_birkhoffAverage_ae}\leanok
  For a finite measure, a measure-preserving \(T\) and an integrable \(g\), the
  Birkhoff averages converge \(\mu\)-a.e.\ to the conditional expectation of
  \(g\) onto the invariant \(\sigma\)-algebra:
  \[
    \texttt{birkhoffAverage}\,\R\,T\,g\,n\,x \to
      \big(\mu[g \mid \texttt{invariants}\,T]\big)(x).
  \]
\end{theorem}
\begin{proof}\leanok
  \uses{measure_setOf_lt_limsup_eq_zero,ae_bddAbove_birkhoffAverage}
  Unioning the null superlevel sets of \cref{measure_setOf_lt_limsup_eq_zero}
  over \(\varepsilon = 1/(k+1)\) gives \(\limsup A_\bullet(g) \le
  \mu[g\mid\mathcal I]\) a.e.; applying the same to \(-g\) (and using
  \(\limsup(-a) = -\liminf a\)) gives \(\mu[g\mid\mathcal I] \le
  \liminf A_\bullet(g)\) a.e. With a.e.\ boundedness
  \cref{ae_bddAbove_birkhoffAverage}, the sandwich
  \(\limsup \le \mu[g\mid\mathcal I] \le \liminf\) forces convergence to
  \(\mu[g\mid\mathcal I]\).
\end{proof}

\begin{corollary}[Ergodic case of Birkhoff]\label{tendsto_birkhoffAverage_ae_integral}
  \lean{Oseledets.tendsto_birkhoffAverage_ae_integral}\leanok
  For an ergodic \(T\) on a probability space and integrable \(g\), the Birkhoff
  averages converge \(\mu\)-a.e.\ to the space average \(\int g\, d\mu\).
\end{corollary}
\begin{proof}\leanok
  \uses{tendsto_birkhoffAverage_ae,condExp_invariants_comp_self}
  By \cref{tendsto_birkhoffAverage_ae} the limit is \(\mu[g\mid\mathcal I]\),
  which is a.e.\ \(T\)-invariant \cref{condExp_invariants_comp_self}; ergodicity
  forces it a.e.\ constant. Its integral equals \(\int g\) (conditional
  expectation preserves the integral) and equals the constant times
  \(\mu(X) = 1\), so the constant is \(\int g\).
\end{proof}

\section{Kingman}

\begin{definition}[Subadditive cocycle]\label{IsSubadditiveCocycle}
  \lean{Oseledets.IsSubadditiveCocycle}\leanok
  A sequence \(g : \N \to X \to \R\) is a \emph{subadditive cocycle} over \(T\)
  when
  \[
    g\,(m+n)\,x \le g\,m\,x + g\,n\,(T^{[m]}x)
      \qquad \text{for all } m, n, x.
  \]
  For \(g_n = \log\norm{A^{(n)}}\) this is submultiplicativity of the operator
  norm composed with the cocycle identity.
\end{definition}

\begin{lemma}[Singleton and block subadditivity]\label{le_birkhoffSum_one}
  \lean{Oseledets.IsSubadditiveCocycle.le_birkhoffSum_one}\leanok
  For a subadditive cocycle and \(n : \N\), one has
  \(g\,(n+1)\,x \le \texttt{birkhoffSum}\,T\,(g\,1)\,(n+1)\,x\); more generally,
  for any decomposition of \([0,N)\) into \(k+1\) consecutive blocks of lengths
  \(\ell_0, \dots, \ell_k\), the cocycle is dominated by the sum of the block
  values along the orbit at the frontiers \(T^{[\sum_{j<i}\ell_j]}x\).
\end{lemma}
\begin{proof}\leanok
  \uses{IsSubadditiveCocycle}
  Both are inductions that peel off the last block and apply the defining
  subadditivity at the split point. The statement is restricted to nonempty
  decompositions (and to index \(n+1\)) because subadditivity at \((0,0)\) only
  forces \(0 \le g\,0\,x\), the wrong sign for a one-sided bound at \(0\).
\end{proof}

\begin{definition}[Normalized cocycle]\label{cdiv}
  \lean{Oseledets.Kingman.cdiv}\leanok
  The \emph{normalized cocycle} is
  \(\texttt{cdiv}\,g\,n\,x = g\,(n+1)\,x / (n+1)\), the sequence whose a.e.\
  limit is the content of Kingman's theorem; its \texttt{EReal} coercion is
  \(\texttt{ecdiv}\,g\,n\,x\), used where \(-\infty\) cannot be excluded a
  priori.
\end{definition}

\begin{lemma}[Fekete limit of the normalized integrals]\label{exists_fekete}
  \lean{Oseledets.Kingman.exists_fekete}\leanok
  For a measure-preserving \(T\), an integrable subadditive cocycle \(g\) whose
  normalized integrals are bounded below, the sequence
  \((\int g\,(n+1)\,d\mu)/(n+1)\) converges to the Fekete constant \(\gamma\).
\end{lemma}
\begin{proof}\leanok
  \uses{IsSubadditiveCocycle}
  Integrating the cocycle inequality and using measure-preservation
  (\(\int g\,n\,(T^{[m]}\cdot) = \int g\,n\)) shows the integral sequence
  \(a_n = \int g\,n\) is subadditive in Fekete's sense. The \((n+1)\)-indexed
  lower bound is bridged by hand to a lower bound on \(a_n/n\) (the \(n=0\) term
  is \(0\)), and Fekete's lemma delivers convergence of \(a_n/n\), hence of the
  shifted sequence, to \(\gamma = \inf_n a_n/n\).
\end{proof}

\begin{lemma}[Invariance from a one-sided orbital bound]\label{ae_eq_comp_of_le_comp}
  \lean{Oseledets.Kingman.ae_eq_comp_of_le_comp}\leanok
  For a finite measure, measure-preserving \(T\) and an a.e.\ measurable \(F\)
  with \(F\,x \le F\,(Tx)\) for a.e.\ \(x\), one has \(F \circ T =^{\text{a.e.}}
  F\).
\end{lemma}
\begin{proof}\leanok
  For each rational \(c\) the upper level set \(\{c \le F\}\) is null-measurable
  (via a measurable representative) and a.e.\ contained in its preimage
  \(T^{-1}\{c \le F\}\), which has equal finite measure; an a.e.\ subset of
  equal finite measure is a.e.\ equal. Ranging over all rational \(c\) and
  collecting the a.e.\ statements gives the invariance.
\end{proof}

\begin{lemma}[Envelopes are a.e.\ measurable and \(T\)-invariant]\label{limsup_div_comp_ae}
  \lean{Oseledets.Kingman.limsup_div_comp_ae}\leanok
  For a finite measure, measure-preserving \(T\), an integrable subadditive
  cocycle with normalized integrals bounded below, the limsup envelope
  \(f_+(x) = \limsup_n \texttt{cdiv}\,g\,n\,x\) (and likewise the liminf
  envelope \(f_-\)) is a.e.\ \(T\)-invariant.
\end{lemma}
\begin{proof}\leanok
  \uses{cdiv,ae_eq_comp_of_le_comp,exists_fekete}
  The subadditive bound \(\texttt{cdiv}\,g\,n\,x \le g\,1\,x/(n+1) + g\,n\,(Tx)/(n+1)\)
  differs from \(\texttt{cdiv}\,g\,n\,(Tx)\) by a null sequence, so the
  vanishing-perturbation lemma for limsup gives the pointwise comparison
  \(f_+(x) \le f_+(Tx)\) wherever the cocycle is a.e.\ bounded (above and
  below) at \(x\) and \(Tx\). Feeding this into
  \cref{ae_eq_comp_of_le_comp} yields a.e.\ invariance.
\end{proof}

\begin{lemma}[Integrability of the limsup envelope]\label{int_limsup_div_integrable}
  \lean{Oseledets.Kingman.int_limsup_div_integrable}\leanok
  Under the same hypotheses, the limsup envelope \(f_+\) is integrable.
\end{lemma}
\begin{proof}\leanok
  \uses{cdiv,tendsto_birkhoffAverage_ae,exists_fekete}
  The nonnegative Fatou defect \(d_n(x) = \texttt{birkhoffAverage}\,\R\,T\,(g\,1)\,
  (n+1)\,x - \texttt{cdiv}\,g\,n\,x \ge 0\) (singleton subadditivity) controls
  the envelope. The \texttt{ENNReal} Fatou inequality \(\int^- \liminf u_n \le
  \liminf \int^- u_n\) applied to \(u_n = \texttt{ofReal}\,d_n\), together with
  \(\int d_n = \int g\,1 - a_{n+1}/(n+1) \to \int g\,1 - \gamma < \infty\)
  (\cref{exists_fekete}) and Birkhoff convergence of \(A_{n+1}(g\,1)\) to a loose
  envelope \(B\) (\cref{tendsto_birkhoffAverage_ae}), shows \(B - f_+ \ge 0\) has
  finite lower integral, hence is integrable; therefore \(f_+ = B - (B - f_+)\)
  is integrable.
\end{proof}

\begin{proposition}[Hard direction: \(\limsup \le \liminf\) a.e.]\label{ae_ereal_limsup_le_liminf}
  \lean{Oseledets.Kingman.ae_ereal_limsup_le_liminf}\leanok
  Under the same hypotheses, for a.e.\ \(x\) the \texttt{EReal} limsup of the
  normalized cocycle is dominated by its liminf.
\end{proposition}
\begin{proof}\leanok
  \uses{cdiv,le_birkhoffSum_one,tendsto_birkhoffAverage_ae,ae_eq_comp_of_le_comp}
  This is the stopping-time / greedy block argument of Katznelson--Weiss and
  Karlsson. After the WLOG shift to the nonpositive process \(\tilde g\,(n+1)
  \le 0\), fix \(\varepsilon, M > 0\) and set \(h = \mu[\max(f_-, -M) \mid
  \mathcal I]\). A greedy two-type partition of \([0,n)\) into ``good'' blocks
  (where the stopping time \(\tau \le L\) realizing \(\tilde g\,\tau \le
  \tau(h+\varepsilon)\) exists) and short singletons (bad/overrun) bounds, via
  block subadditivity \cref{le_birkhoffSum_one},
  \(\tilde g\,n\,x/n \le (h\,x+\varepsilon)(1 - (L-1)/n -
  \texttt{birkhoffAverage}\,\mathbf 1_{B_L})\). Letting \(n \to \infty\)
  (Birkhoff, \cref{tendsto_birkhoffAverage_ae}), then \(L \to \infty\) (the bad
  sets \(B_L\) shrink to a null set), \(M \to \infty\) and \(\varepsilon \to 0\)
  yields \(f_+ \le f_-\) a.e. The whole argument is carried in \texttt{EReal} to
  keep the bookkeeping clean near \(-\infty\).
\end{proof}

\begin{theorem}[Kingman core: a.e.\ existence of an integrable limit]\label{ae_tendsto_cdiv}
  \lean{Oseledets.Kingman.ae_tendsto_cdiv}\leanok
  For a finite measure, measure-preserving measurable \(T\), an integrable
  subadditive cocycle with normalized integrals bounded below, there is an
  integrable \(G\) with \(\texttt{cdiv}\,g\,n\,x \to G\,x\) for \(\mu\)-a.e.\
  \(x\).
\end{theorem}
\begin{proof}\leanok
  \uses{cdiv,int_limsup_div_integrable,ae_ereal_limsup_le_liminf}
  Take \(G = f_+\), integrable by \cref{int_limsup_div_integrable}. On the
  a.e.\ good set the \texttt{EReal} limsup \(e\) satisfies \(\bot < e \le B <
  \top\) (finiteness from the loose envelope and the Fatou step), and
  \(\liminf = \limsup = e\) \cref{ae_ereal_limsup_le_liminf}; in a complete
  linear order equal liminf and limsup force convergence to \(e\). Transferring
  to \(\R\) gives \(\texttt{cdiv}\,g\,n\,x \to e^{\text{toReal}} = f_+\,x\).
\end{proof}

\begin{theorem}[Kingman's subadditive ergodic theorem]\label{tendsto_kingman}
  \lean{Oseledets.tendsto_kingman}\leanok
  For a finite measure, a measure-preserving \(T\), an integrable subadditive
  cocycle \(g\) whose normalized integrals are bounded below, there is a
  \(T\)-invariant integrable \(G\) with
  \[
    n^{-1}\, g\,n\,x \to G\,x
      \qquad \text{for } \mu\text{-a.e.\ } x.
  \]
\end{theorem}
\begin{proof}\leanok
  \uses{ae_tendsto_cdiv,limsup_div_comp_ae,int_limsup_div_integrable,isSubadditiveCocycle_logNorm,lambdaBar}
  Take \(G = f_-\). On the a.e.\ set where the normalized cocycle is bounded,
  \(\liminf \le \limsup\) is trivial and \(\limsup \le \liminf\) is the hard
  direction inside \cref{ae_tendsto_cdiv}, so \(f_- =^{\text{a.e.}} f_+\); the
  sandwich \(\limsup \le f_- \le \liminf\) yields pointwise convergence to
  \(f_-\), and reindexing removes the \(n=0\) term. Invariance is
  \cref{limsup_div_comp_ae} (liminf variant), and integrability follows from
  \(f_- =^{\text{a.e.}} f_+\) with \cref{int_limsup_div_integrable}.
\end{proof}

\begin{corollary}[Kingman, ergodic case]\label{tendsto_kingman_ergodic}
  \lean{Oseledets.tendsto_kingman_ergodic}\leanok
  For an ergodic \(T\) on a probability space, an integrable subadditive cocycle
  with normalized integrals bounded below, there is a constant \(c\) with
  \(n^{-1}\,g\,n\,x \to c\) for \(\mu\)-a.e.\ \(x\).
\end{corollary}
\begin{proof}\leanok
  \uses{tendsto_kingman}
  Kingman's theorem \cref{tendsto_kingman} gives a \(T\)-invariant integrable
  limit \(G\); ergodicity forces an a.e.\ \(T\)-invariant integrable function to
  be a.e.\ constant, so \(G =^{\text{a.e.}} c\) and the limit is \(c\). (That
  constant is the Fekete infimum \(\gamma\); only a.e.-constancy is asserted, as
  this is what the multiplicative ergodic theorem consumes.)
\end{proof}
